
\section{\BookRecordIndex{T0257}粘合映射的导数}\label{sec:10-6}

本节先给出一些重要的技术结果，它们将用于证明命题~\ref{prop:10-5-4}；
上面正是用该命题建立了\BookRecordIndex{T0256}粘合定理的 (i)--(iii)。本节最后证明关于\BookRecordIndex{T0257}粘合映射
保持定向的 (iv)。这些论证虽然有时较长，但都是直接的。以下始终固定两个
同调类 $A^0,A^\infty\in H_2(M;\mathbb Z)$，以及一个\BookRecordIndex{T0088}殆复结构\BookRecordIndex{T0406}正则对
\[
  (J^0,J^\infty)\in
  \mathcal J_{\mathrm{reg}}(S^2;A^{0,\infty}).
\]
与前面一样，简记
\[
  \mathcal M(c):=\mathcal M(A^{0,\infty};J^{0,\infty},c),
\]
并常把 $\mathcal M(c)$ 的元素 $(u^0,u^\infty)$ 简记为 $u$。此外，把乘积
$\delta R$ 记作 $r$。于是，$(\delta,R)\in\mathcal A(\delta_0)$ 意味着
$r>1/\delta_0$。

第一个引理为中间映射族 $u^{0,r},u^{\infty,r}$ 上线性化\BookRecordIndex{T0210}柯西--黎曼算子的
\BookRecordIndex{T0392}右逆 $Q_{0,\infty,r}$ 建立精确估计，特别是证明上面用到的
估计~\eqref{eq:10-5-2}。这里必须谨慎，因为我们只对 $u^{a,r}$ 的一阶导数
有一致界，对其高阶导数则没有一致界。因此，只能使用沿这些曲线的向量场之
$L^p$ 与 $W^{1,p}$ 范数。

\begin{lemma}\label{lem:10-6-1}
固定 $p>2$ 和 $c>0$。则存在正数 $\delta_0,c_0$，使得对所有
$(u^0,u^\infty)\in\mathcal M(c)$、所有
$(\delta,R)\in\mathcal A(\delta_0)$，令 $r:=\delta R$ 后均有：

\begin{enumerate}[label=(\roman*)]
\item 对每个
$\xi=(\xi^0,\xi^\infty)\in W_{u^{0,\infty},r}^{1,p}$，
\[
  \|\xi\|_{W^{1,p}}
  \le c_0\bigl(\|D_{0,\infty,r}\xi\|_{L^p}
                 +\|\xi\|_{L^p}\bigr).
\]

\item 对每个
$\xi=(\xi^0,\xi^\infty)\in W_{u^{0,\infty},r}^{1,p}$，
\[
  D_{0,\infty,r}\xi=0
  \quad\Longrightarrow\quad
  \|\xi\|_{W^{1,p}}\le c_0\|\xi\|_{L^2}.
\]

\item 对每个
\[
  \eta=(\eta^0,\eta^\infty)
  \in L_{u^{0,r},J^0}^p\times L_{u^{\infty,r},J^\infty}^p
\]
都有
\[
  \|Q_{0,\infty,r}\eta\|_{W^{1,p}}
  \le c_0\|\eta\|_{L^p}.
\]
\end{enumerate}
\end{lemma}

\begin{proof}
当 $r$ 充分大时，经由\BookRecordIndex{T0263}平行移动等同定义域和靶之后，算子
$D_{0,\infty,r}$ 在算子范数意义下任意接近原算子 $D_{0,\infty}$。
所以 (i) 立即由 $D_{0,\infty}$ 的相应估计\BookRecordIndex{T0326}推出。为了证明其余估计，考虑
下面的抽象泛函分析情形。设
\[
  D:\mathcal W\longrightarrow\mathcal L
\]
是两个巴拿赫空间之间指标为 $d$ 的满射\BookRecordIndex{T0128}弗雷德霍姆算子。（可把
$\mathcal W=W^{1,p}$、$\mathcal L=L^p$、$D=D_{0,\infty}$
是一阶椭圆算子的情形作为模型。）设 $\mathcal W$ 配备内积
$\langle\cdot,\cdot\rangle$，并把相应范数记作
\[
  |\xi|:=\sqrt{\langle\xi,\xi\rangle}.
\]
（可把它想成 $L^2$ 内积。）再假设存在正数 $c,c_D$，使每个
$\xi\in\mathcal W$ 都有
\begin{equation}\label{eq:10-6-1}
  |\xi|\le c\|\xi\|_{\mathcal W},
\end{equation}
并且
\begin{equation}\label{eq:10-6-2}
  D\xi=0\quad\Longrightarrow\quad
  \|\xi\|_{\mathcal W}\le c_D|\xi|.
\end{equation}
在当前情形，后一条件成立是因为 $D$ 的核有限维。记
$Q:\mathcal L\to\mathcal W$ 为
$D$ 的\BookRecordIndex{T0392}右逆，其像是关于上述内积的 $\ker D$ 之正交补。

现设 $D':\mathcal W\to\mathcal L$ 是另一有界线性算子
（例如 $D_{0,\infty,r}$），且
\[
  \varepsilon:=\|D'-D\|\,\|Q\|<1.
\]
由于 $DQ=\operatorname{id}$，有
$\|D'Q-\operatorname{id}\|<1$，故 $D'$ 满射并有\BookRecordIndex{T0392}右逆
$Q(D'Q)^{-1}$。不过，我们要考察的是\BookRecordIndex{T0392}右逆
$Q':\mathcal L\to\mathcal W$，其像为
$\ker D'$ 的正交补。

先注意，若 $D'\zeta=0$，则
$\|QD\zeta\|_{\mathcal W}\le\varepsilon\|\zeta\|_{\mathcal W}$，且
$\|\zeta-QD\zeta\|_{\mathcal W}
\ge(1-\varepsilon)\|\zeta\|_{\mathcal W}$。因而
\begin{equation}\label{eq:10-6-3}
  D'\zeta=0
  \quad\Longrightarrow\quad
  \|QD\zeta\|_{\mathcal W}
  \le\frac{\varepsilon}{1-\varepsilon}
      \|\zeta-QD\zeta\|_{\mathcal W}.
\end{equation}
由于 $\zeta-QD\zeta\in\ker D$，有
\[
  \|\zeta-QD\zeta\|_{\mathcal W}
  \le c_D|\zeta-QD\zeta|
  \le c_D|\zeta|.
\]
最后一个不等式成立，是因为 $\zeta-QD\zeta$ 是 $\zeta$ 到 $\ker D$
上的正交投影。另一方面，上面已经看到
$\|(1-\varepsilon)\zeta\|_{\mathcal W}
\le\|\zeta-QD\zeta\|_{\mathcal W}$。所以
\[
  D'\zeta=0
  \quad\Longrightarrow\quad
  \|\zeta\|_{\mathcal W}\le\frac{c_D}{1-\varepsilon}|\zeta|.
\]
这就证明了 (ii)。

现设 $\xi\in\mathcal W$ 与 $\ker D'$ 正交。取 $\ker D$ 的一个标准正交基
$e_1,\ldots,e_d$，并定义 $\ker D'$ 的基
$e'_1,\ldots,e'_d$，使
\[
  e'_i-QDe'_i=e_i,
  \qquad i=1,\ldots,d.
\]
（映射
$\ker D'\to\ker D:\zeta\mapsto\zeta-QD\zeta$
是两个核之间的同构。）由于 $\xi-QD\xi\in\ker D$，有
\[
  \xi-QD\xi=\sum_{i=1}^d\langle\xi,e_i\rangle e_i.
\]
此外，
\[
  \langle\xi,e_i\rangle
  =\langle\xi,e_i-e'_i\rangle
  =\langle\xi,-QDe'_i\rangle;
\]
所以由式~\eqref{eq:10-6-3}，并结合
式~\eqref{eq:10-6-1}、\eqref{eq:10-6-2}，得到
\[
  |\langle\xi,e_i\rangle|
  \le c\|QDe'_i\|_{\mathcal W}\|\xi\|_{\mathcal W}
  \le\frac{c\varepsilon}{1-\varepsilon}
      \|e_i\|_{\mathcal W}\|\xi\|_{\mathcal W}
  \le\frac{cc_D\varepsilon}{1-\varepsilon}\|\xi\|_{\mathcal W}.
\]
结合前一恒等式，得
\[
  \|\xi-QD\xi\|_{\mathcal W}
  \le c_D|\xi-QD\xi|
  =c_D\left(\sum_{i=1}^d\langle\xi,e_i\rangle^2\right)^{1/2}
  \le\frac{\sqrt d\,c c_D^2\varepsilon}{1-\varepsilon}
      \|\xi\|_{\mathcal W}.
\]
因而
\[
\begin{aligned}
  \|\xi\|_{\mathcal W}
  &\le\|QD\xi\|_{\mathcal W}+\|\xi-QD\xi\|_{\mathcal W}\\
  &\le\|Q\|\,\|D'\xi\|_{\mathcal L}
     +\|Q\|\,\|D'-D\|\,\|\xi\|_{\mathcal W}
     +\|\xi-QD\xi\|_{\mathcal W}\\
  &\le\|Q\|\,\|D'\xi\|_{\mathcal L}
     +\varepsilon\left(1+
       \frac{\sqrt d\,c c_D^2}{1-\varepsilon}\right)
       \|\xi\|_{\mathcal W}.
\end{aligned}
\]
若
\[
  \varepsilon\le\frac12,
  \qquad
  \varepsilon(1+2\sqrt d\,c c_D^2)\le\frac12,
\]
就\BookRecordIndex{T0326}推出每个与 $\ker D'$ 正交的 $\xi\in\mathcal W$ 都满足
\[
  \|\xi\|_{\mathcal W}\le2\|Q\|\,\|D'\xi\|_{\mathcal L}.
\]
回忆 $\|D'-D\|=\varepsilon/\|Q\|$，且 $D'$ 满射，故 $D'\xi$
遍历 $\mathcal L$ 的所有元素。于是存在 $\delta>0$，使
\[
  \|D'-D\|<\delta\quad\Longrightarrow\quad\|Q'\|\le2\|Q\|.
\]
$\delta$ 要取得多小，只依赖于 $Q$ 的算子范数以及常数 $d,c,c_D$。
这就证明了引理~\ref{lem:10-6-1}。
\end{proof}

接着考察近似逆 $T_{u^R}$ 随 $u$ 沿\BookRecordIndex{T0246}模空间 $\mathcal M(c)$ 中的道路
$u_t:=(u_t^0,u_t^\infty)$ 变化时的协变导数。为简化计算，对向量场使用
\BookRecordIndex{T0034}Levi--Civita 联络 $\nabla$；但对 $1$-形式必须使用\BookRecordIndex{T0052}埃尔米特联络
$\widetilde\nabla$，因为只有后者保持 $(0,1)$-形式。回忆
$\mathcal A(\delta_0)$ 表示容许偶对 $(\delta,R)$ 的集合，即
$(\delta,R)\in\mathcal A(\delta_0)$ 当且仅当
$0<\delta<\delta_0$ 且 $\delta R>1/\delta_0$。

\begin{lemma}\label{lem:10-6-2}
固定 $c>0$ 和 $p>2$。则存在正数 $\delta_0,c_0$，具有下述性质。
固定 $(\delta,R)\in\mathcal A(\delta_0)$；设
\[
  \mathbb R\longrightarrow\mathcal M(c):
  t\longmapsto(u_t^0,u_t^\infty)
\]
是一条光滑道路，并设 $t\mapsto u_t^R$ 是由
式~\eqref{eq:10-2-2} 定义的相应\BookRecordIndex{T0395}预粘合曲线道路。再设
$t\mapsto\eta_t^R$ 是沿 $u_t^R$ 的一族光滑 $(0,1)$-形式，即每个 $t$ 都有
\[
  \eta_t^R\in\Omega^{0,1}(S^2,(u_t^R)^*TM).
\]
则
\[
  \bigl\|\nabla_t(T_{u_t^R}\eta_t^R)
    -T_{u_t^R}\widetilde\nabla_t\eta_t^R\bigr\|_{1,p,R}
  \le c_0\|\partial_tu_t^R\|_{1,\infty,R}
         \|\eta_t^R\|_{0,p,R}.
\]
\end{lemma}

\begin{proof}
先在一个没有\BookRecordIndex{T0394}预粘合映射的简单情形下证明类似估计。设有一族球面映射
\[
  \mathbb R\times S^2\longrightarrow M:(t,z)\longmapsto u_t(z),
\]
记相应的线性化\BookRecordIndex{T0210}柯西--黎曼算子族为 $D_{u_t}$，并设
$\eta_t\in\Omega^{0,1}(S^2,u_t^*TM)$ 是沿 $u_t$ 的一族光滑
$(0,1)$-形式。始终假设存在常数 $c$，使每个 $t$ 都有
$\|du_t\|_{L^\infty}\le c$。还假设每个 $D_t:=D_{u_t}$ 都满射，
并记 $Q_t:=Q_{u_t}$ 为其像等于核之 $L^2$-正交补的\BookRecordIndex{T0392}右逆。定义
$\xi_t\in\Omega^0(S^2,u_t^*TM)$ 为
\[
  \xi_t:=Q_t\eta_t,
  \qquad \eta_t=D_t\xi_t.
\]
目标是得到估计
\begin{equation}\label{eq:10-6-4}
  \|\nabla_t\xi_t-Q_t\widetilde\nabla_t\eta_t\|_{W^{1,p}}
  \le C\|\partial_tu_t\|_{W^{1,\infty}}\|\eta_t\|_{L^p}.
\end{equation}
我们将使用不等式
\begin{equation}\label{eq:10-6-5}
  D_t\xi=0\quad\Longrightarrow\quad
  \|\xi\|_{W^{1,p}}\le c_1\|\xi\|_{L^2},
\end{equation}
\begin{equation}\label{eq:10-6-6}
  \|Q_t\eta\|_{W^{1,p}}\le c_2\|\eta\|_{L^p},
\end{equation}
\begin{equation}\label{eq:10-6-7}
  \|\xi\|_{L^2}\le c_3\|\xi\|_{W^{1,p}}.
\end{equation}
还要使用交换子估计
\begin{equation}\label{eq:10-6-8}
  \|\widetilde\nabla_t(D_t\xi_t)-D_t\nabla_t\xi_t\|_{L^p}
  \le c_4\|\partial_tu_t\|_{W^{1,\infty}}
          \|\xi_t\|_{W^{1,p}}.
\end{equation}
为证明式~\eqref{eq:10-6-8}，对 $1$-形式场
$t\mapsto D_t\xi_t$ 作协变微分，得到
\[
  \widetilde\nabla_t(D_t\xi_t)-D_t\nabla_t\xi_t
  =B_1(\nabla_tdu_t,\xi_t)
   +B_2(\partial_tu_t,\nabla\xi_t)
   +B_3(\partial_tu_t,du_t,\xi_t).
\]
这里 $B_1,B_2$ 是取值于 $TM$ 的 $TM$ 上双线性形式，$B_3$ 是取值于
$TM$ 的 $TM$ 上三线性形式；这些多重线性形式均一致有界。逐项考察右端即得
式~\eqref{eq:10-6-8}。

下面说明最后四个不等式如何合并\BookRecordIndex{T0326}推出式~\eqref{eq:10-6-4}。提醒读者，
$D_t$ 表示\BookRecordIndex{T0210}柯西--黎曼算子族；对 $t$ 的微分可记为 $\nabla_t$（对 $M$ 中张量
作微分时），也可记为 $\partial_t$（对以 $t$ 为参数的道路上的函数作微分时）。
首先，由式~\eqref{eq:10-6-6} 和~\eqref{eq:10-6-8}，
\[
\begin{aligned}
  \|Q_tD_t\nabla_t\xi_t-Q_t\widetilde\nabla_t\eta_t\|_{W^{1,p}}
  &\le c_2\|D_t\nabla_t\xi_t
              -\widetilde\nabla_t(D_t\xi_t)\|_{L^p}\\
  &\le c_2c_4\|\partial_tu_t\|_{W^{1,\infty}}
                 \|\xi_t\|_{W^{1,p}}\\
  &\le c_2^2c_4\|\partial_tu_t\|_{W^{1,\infty}}
                  \|\eta_t\|_{L^p}.
\end{aligned}
\]
现设 $t\mapsto\zeta_t$ 是 $\ker D_t$ 中的一条道路，且
$\|\zeta_t\|_{L^2}=1$。由于 $D_t\zeta_t=0$，有
\[
\begin{aligned}
  \|Q_tD_t\nabla_t\zeta_t\|_{L^p}
  &\le c_2\|D_t\nabla_t\zeta_t
              -\widetilde\nabla_t(D_t\zeta_t)\|_{L^p}\\
  &\le c_2c_4\|\partial_tu_t\|_{W^{1,\infty}}
                 \|\zeta_t\|_{W^{1,p}}\\
  &\le c_1c_2c_4\|\partial_tu_t\|_{W^{1,\infty}}.
\end{aligned}
\]
此外，
\[
  0=\partial_t\langle\xi_t,\zeta_t\rangle
   =\langle\nabla_t\xi_t,\zeta_t\rangle
    +\langle\xi_t,\nabla_t\zeta_t\rangle,
\]
故
\[
\begin{aligned}
  |\langle\nabla_t\xi_t,\zeta_t\rangle|
  &=|\langle\xi_t,\nabla_t\zeta_t\rangle|\\
  &=|\langle\xi_t,Q_tD_t\nabla_t\zeta_t\rangle|\\
  &\le c_1c_2c_3c_4
       \|\partial_tu_t\|_{W^{1,\infty}}\|\xi_t\|_{W^{1,p}}.
\end{aligned}
\]
取 $\ker D_t$ 的一个标准正交\BookRecordIndex{T0067}标架
$\zeta_{1t},\ldots,\zeta_{mt}$，其中 $m$ 是 $D$ 的\BookRecordIndex{T0129}弗雷德霍姆指标。于是
\[
\begin{aligned}
  \|\nabla_t\xi_t-Q_tD_t\nabla_t\xi_t\|_{W^{1,p}}
  &\le c_1\|\nabla_t\xi_t-Q_tD_t\nabla_t\xi_t\|_{L^2}\\
  &=c_1\left(
     \sum_{i=1}^m\langle\nabla_t\xi_t,\zeta_{it}\rangle^2
     \right)^{1/2}\\
  &\le\sqrt m\,c_1^2c_2c_3c_4
       \|\partial_tu_t\|_{W^{1,\infty}}\|\xi_t\|_{W^{1,p}}\\
  &\le\sqrt m\,c_1^2c_2^2c_3c_4
       \|\partial_tu_t\|_{W^{1,\infty}}\|\eta_t\|_{L^p}.
\end{aligned}
\]
这便以 $C:=c_2^2c_4(1+\sqrt m\,c_1^2c_3)$ 证明了
式~\eqref{eq:10-6-4}。

现在必须把式~\eqref{eq:10-6-4} 用于如下情形：$u_t$ 换成由
$u_t^0,u_t^\infty$ 按注~\ref{rem:10-2-2} 构造的\BookRecordIndex{T0226}连通对
$(u_t^{0,r},u_t^{\infty,r})$，而我们所关心的是\BookRecordIndex{T0395}预粘合曲线 $u_t^R$ 上的
算子 $D,T$。于是，令 $t\mapsto\eta_t^R$ 是引理假设中沿\BookRecordIndex{T0394}预粘合
$u_t^R$ 的 $(0,1)$-形式族，记
\[
  \eta_t:=(\eta_t^0,\eta_t^\infty)
\]
为命题~\ref{prop:10-5-1} 的证明中定义的、沿
$u_t^{0,r},u_t^{\infty,r}$ 的相应 $(0,1)$-形式偶对，并简记
\[
  D_{t,r}:=D_{u_t^{0,r},u_t^{\infty,r}},
  \qquad
  Q_{t,r}:=Q_{u_t^{0,r},u_t^{\infty,r}}.
\]
令
\[
  \xi_t:=(\xi_t^0,\xi_t^\infty)
  :=Q_{t,r}(\eta_t^0,\eta_t^\infty),
  \qquad
  \xi_t^R:=T_{u_t^R}\eta_t^R.
\]
这样，$\xi_t^R$ 是由命题~\ref{prop:10-5-1} 的证明中
式~\eqref{eq:10-5-4} 从偶对 $\xi_t$ 得到的。

由引理~\ref{lem:10-6-1}，算子 $D_{t,r},Q_{t,r}$ 以只依赖 $c,p$ 的
适当常数 $c_1,c_2,c_3,c_4$ 满足不等式~\eqref{eq:10-6-5}--
\eqref{eq:10-6-8}。所以，由刚才证明的结论，
\begin{equation}\label{eq:10-6-9}
  \|\nabla_t\xi_t-Q_{t,r}\widetilde\nabla_t\eta_t\|_{W^{1,p}}
  \le C\|\partial_tu_t^R\|_{1,\infty,R}
         \|\eta_t^R\|_{0,p,R}.
\end{equation}
这里还使用了不等式
\[
  \|\eta_t\|_{L^p}\le2\|\eta_t^R\|_{0,p,R},
  \qquad
  \|\partial_tu_t^{0,r}\|_{W^{1,\infty}}
  +\|\partial_tu_t^{\infty,r}\|_{W^{1,\infty}}
  \le2\|\partial_tu_t^R\|_{1,\infty,R}.
\]
式~\eqref{eq:10-5-4} 的右端定义有界线性算子
\[
  W_{u^{0,\infty},r}^{1,p}\longrightarrow W_{u^R}^{1,p}:
  \xi=(\xi^0,\xi^\infty)\longmapsto\xi^R,
\]
而且它与沿每条道路
$t\mapsto u_t:=(u_t^0,u_t^\infty)$ 的协变导数 $\nabla_t$ 交换。
（这是因为\BookRecordIndex{T0188}截断函数 $\beta$ 只依赖定义域中的点 $z$，与 $t$ 无关。）因此，
由式~\eqref{eq:10-6-9}，存在常数 $C'>0$，使
\[
  \|\nabla_t\xi_t^R-T_{u_t^R}\widetilde\nabla_t\eta_t^R\|_{1,p,R}
  \le C'\|\partial_tu_t^R\|_{1,\infty,R}
          \|\eta_t^R\|_{0,p,R}.
\]
这就证明了引理~\ref{lem:10-6-2}。
\end{proof}

下一个引理描述 $\mathcal M(c)$ 的切空间与\BookRecordIndex{T0395}预粘合曲线\BookRecordIndex{T0246}模空间的切空间之间的
同构。我们对 $T_u\mathcal M(c)=\ker D_u$ 中的元素使用 $L^2$-范数；
这只是为了方便：当 $u=(u^0,u^\infty)$ 遍历紧空间 $\mathcal M(c)$ 时，
这个有限维空间上的所有范数都以一致常数等价。另一方面，这一范数也很自然，
因为我们正是用 $\ker D_u$ 的 $L^2$-正交补选取唯一\BookRecordIndex{T0392}右逆。为与本节其余记号
一致，以下通常把 $D_{0,\infty}$ 写成 $D_u$。回忆
\[
  Q_{u^R}:=T_{u^R}(D_{u^R}T_{u^R})^{-1}
\]
是 $D_{u^R}$ 的\BookRecordIndex{T0392}右逆，其像等于命题~\ref{prop:10-5-1} 中算子
$T_{u^R}$ 的像。

\begin{lemma}\label{lem:10-6-3}
对每个 $c>0$ 和 $p>2$，存在正数 $\delta_0,c_0$，使得若
$u=(u^0,u^\infty)\in\mathcal M(c)$ 且
$(\delta,R)\in\mathcal A(\delta_0)$，则对每个
$(\xi^0,\xi^\infty)\in T_u\mathcal M(c)$ 都有
\begin{equation}\label{eq:10-6-10}
  \|(\xi^0,\xi^\infty)\|_{L^2}
  \le c_0\bigl\|
    (\operatorname{id}-Q_{u^R}D_{u^R})
    df^R(u)(\xi^0,\xi^\infty)
  \bigr\|_{1,p,R}.
\end{equation}
因而算子
\[
  (\operatorname{id}-Q_{u^R}D_{u^R})df^R(u)
\]
是从 $D_u:=D_{u^0,u^\infty}$ 的核到 $D_{u^R}$ 的核的同构。
\end{lemma}

\begin{proof}
简记
\[
  D^R:=D_{u^R},\qquad Q^R:=Q_{u^R},\qquad T^R:=T_{u^R}.
\]
我们将证明：对
$u=(u^0,u^\infty)\in\mathcal M(c)$ 和
$\xi=(\xi^0,\xi^\infty)\in\ker D_u=T_u\mathcal M(c)$，有
\begin{equation}\label{eq:10-6-11}
  \|D^Rdf^R(u)\xi\|_{0,p,R}
  \le c_2(\delta R)^{-2/p}\|\xi\|_{L^2}.
\end{equation}
把它与式~\eqref{eq:10-5-8} 及式~\eqref{eq:10-5-15} 中的估计
$\|\xi\|_{L^2}\le c_1\|df^R(u)\xi\|_{1,p,R}$ 合并，得到
\[
\begin{aligned}
  c_1\|(\operatorname{id}-Q^RD^R)df^R(u)\xi\|_{1,p,R}
  &\ge c_1\|df^R(u)\xi\|_{1,p,R}
     -c_0c_1\|D^Rdf^R(u)\xi\|_{0,p,R}\\
  &\ge\bigl(1-c_0c_1c_2(\delta R)^{-2/p}\bigr)
     \|\xi\|_{L^2}.
\end{aligned}
\]
所以，只要 $\delta R$ 充分大，式~\eqref{eq:10-6-10} 就成立。

为证明式~\eqref{eq:10-6-11}，简记
$\xi^R:=df^R(u)\xi$。按构造，$D^R\xi^R$ 在\BookRecordIndex{T0170}环域
\[
  \frac1{\delta R}\le|z|\le\frac2{\delta R},
  \qquad
  \frac\delta{2R}\le|z|\le\frac\delta R
\]
之外为零。为了在第一个\BookRecordIndex{T0170}环域 $1/r\le|z|\le2/r$（照例
$r=\delta R$）上估计 $D^R\xi^R$，沿 $\mathcal M(c)$ 中的一条道路对
$u^R$ 的定义公式~\eqref{eq:10-2-2} 求导，以得到 $\xi^R$ 的公式。把道路记作
$t\mapsto u_t:=(u_t^0,u_t^\infty)$，并在 $1/r\le|z|\le2/r$ 上写成
\[
  u_t^0(z)=:\exp_{x_t}(\zeta_t^0(z)),
  \qquad x_t:=u_t^0(0).
\]
于是 $u_t^R:=f^R(u_t^0,u_t^\infty)$ 在该\BookRecordIndex{T0170}环域上由
\[
  u_t^R(z)=\exp_{x_t}(\rho(rz)\zeta_t^0(z))
\]
给出。对这两个恒等式关于 $t$ 求导，并使用注~\ref{rem:10-5-5} 的记号，得
\[
\begin{aligned}
  \partial_tu_t^R
  & =E_1(x_t,\rho(rz)\zeta_t^0(z))\partial_tx_t
    +\rho(rz)E_2(x_t,\rho(rz)\zeta_t^0(z))
      \nabla_t\zeta_t^0(z),\\
  \partial_tu_t^0
  & =E_1(x_t,\zeta_t^0(z))\partial_tu_t^0(0)
    +E_2(x_t,\zeta_t^0(z))\nabla_t\zeta_t^0(z).
\end{aligned}
\]
由第二式解出 $\nabla_t\zeta_t^0(z)$ 并代入第一式，再令 $t=0$，使用
$\partial_tu_t^0|_{t=0}=\xi^0$ 与
$\partial_tu_t^R|_{t=0}=\xi^R$，得到
\begin{equation}\label{eq:10-6-12}
  \xi^R(z)
  =E_1(x,\rho(rz)\zeta^0(z))\xi^0(0)
   +E_2(x,\rho(rz)\zeta^0(z))\widehat\xi^0(z),
\end{equation}
其中 $x:=u^0(0)$，而 $\zeta^0(z),\widehat\xi^0(z)\in T_xM$ 由
\[
  u^0(z)=:\exp_x(\zeta^0(z))
\]
及
\begin{equation}\label{eq:10-6-13}
  \widehat\xi^0(z)
  :=\rho(rz)E_2(x,\zeta^0(z))^{-1}
    \bigl(\xi^0(z)-E_1(x,\zeta^0(z))\xi^0(0)\bigr)
\end{equation}
定义，且 $1/r\le|z|\le2/r$。这个显式公式表明，$\xi^R$ 在该\BookRecordIndex{T0170}环域中的
$C^1$-范数不超过 $\xi^0$ 的 $C^1$-范数的一个常数倍；另一个\BookRecordIndex{T0170}环域
$\delta/(2R)\le|z|\le\delta/R$ 同理。与引理~\ref{lem:10-3-2} 的证明
一样，现利用两个\BookRecordIndex{T0170}环域关于依赖 $R$ 的度量之面积均不超过
$(\delta R)^{-2}$ 的一个常数倍这一事实，\BookRecordIndex{T0326}推出
\[
  \|D^R\xi^R\|_{0,p,R}
  \le c_4(\delta R)^{-2/p}\|(\xi^0,\xi^\infty)\|_{C^1}
  \le c_5(\delta R)^{-2/p}\|(\xi^0,\xi^\infty)\|_{L^2}.
\]
最后一个不等式成立，是因为对 $D_{0,\infty}$ 核中的每个元素，其
$L^2$-范数控制 $C^1$-范数。这就证明了引理~\ref{lem:10-6-3}。
\end{proof}

\begin{exercise}\label{ex:10-6-4}
证明引理~\ref{lem:10-6-3} 的同构
\[
  (\operatorname{id}-Q_{u^R}D_{u^R})df^R(u):
  \ker D_u\longrightarrow\ker D_{u^R}
\]
保持定向。提示：先假设 $D_u,D_{u^R}$ 都是复线性的，并估计该同构的复反线性
部分；再通过同伦论证把一般情形化归为复线性情形。
\end{exercise}

\begin{proof}[命题~\ref{prop:10-5-4} 的证明]
设 $t\mapsto u_t:=(u_t^0,u_t^\infty)$ 是 $\mathcal M(c)$ 中的光滑道路，
记 $u_t^R:=f^R(u_t)$，并取
$\widetilde\xi_t^R\in\operatorname{im}Q_{u_t^R}$，使
\[
  \widetilde u_t^R:=\exp_{u_t^R}(\widetilde\xi_t^R)
  =\iota_c^R(u_t).
\]
以下简记
\[
  \widetilde D_t^R:=D_{\widetilde u_t^R},
  \qquad D_t^R:=D_{u_t^R},
  \qquad T_t^R:=T_{u_t^R},
  \qquad Q_t^R:=Q_{u_t^R}.
\]
由于 $\widetilde\xi_t^R\in\operatorname{im}T_t^R
=\operatorname{im}Q_t^R$，有
\[
  \widetilde\xi_t^R=T_t^R\eta_t^R,
  \qquad
  \eta_t^R:=(D_t^RT_t^R)^{-1}D_t^R\widetilde\xi_t^R.
\]
由式~\eqref{eq:10-5-11} 和~\eqref{eq:10-5-1}，存在常数 $c_2>0$，使
\begin{equation}\label{eq:10-6-14}
  \|\widetilde\xi_t^R\|_{1,p,R}+\|\eta_t^R\|_{0,p,R}
  \le c_2(\delta R)^{-2/p}.
\end{equation}
此外，把式~\eqref{eq:10-6-12} 和~\eqref{eq:10-6-13} 中的
$\xi^0,\xi^\infty,\xi^R$ 分别换成
$\partial_tu_t^0,\partial_tu_t^\infty,\partial_tu_t^R$，可知
$\partial_tu_t^R$ 的 $(1,\infty,R)$-范数可以由偶对
$(\partial_tu_t^0,\partial_tu_t^\infty)$ 的 $C^1$-范数估计。
由于该偶对属于 $D_{u_t^0,u_t^\infty}$ 的核，其 $L^2$-范数控制
$C^1$-范数。因此存在常数 $c_3>0$，使
\begin{equation}\label{eq:10-6-15}
  \|\partial_tu_t^R\|_{1,\infty,R}
  \le c_3\left(
    \|\partial_tu_t^0\|_{L^2}
    +\|\partial_tu_t^\infty\|_{L^2}
  \right).
\end{equation}
又由式~\eqref{eq:10-2-2} 中 $u_t^R$ 的定义，$(0,1)$-形式
$D_t^R\partial_tu_t^R$ 支撑在\BookRecordIndex{T0170}环域
\[
  \frac1{\delta R}\le|z|\le\frac2{\delta R},
  \qquad
  \frac\delta{2R}\le|z|\le\frac\delta R
\]
中；两个\BookRecordIndex{T0170}环域的面积都不超过 $4\pi(\delta R)^{-2}$。所以存在常数
$c_4>0$，使
\begin{equation}\label{eq:10-6-16}
  \|D_t^R\partial_tu_t^R\|_{0,p,R}
  \le c_4(\delta R)^{-2/p}\|\partial_tu_t^R\|_{1,\infty,R}.
\end{equation}
由引理~\ref{lem:10-6-2}，还存在常数 $c_5>0$，使
\begin{equation}\label{eq:10-6-17}
  \|\nabla_t\widetilde\xi_t^R
     -T_t^R\widetilde\nabla_t\eta_t^R\|_{1,p,R}
  \le c_5\|\partial_tu_t^R\|_{1,\infty,R}
          \|\eta_t^R\|_{0,p,R}.
\end{equation}

下面说明最后四个不等式如何\BookRecordIndex{T0326}推出命题结论。对恒等式
\[
  \bar\partial_{J^R}
    \bigl(\exp_{u_t^R}(\widetilde\xi_t^R)\bigr)=0
\]
关于 $t$ 求导，得
\begin{equation}\label{eq:10-6-18}
  \widetilde D_t^R
  \bigl(E_1^R\partial_tu_t^R
        +E_2^R\nabla_t\widetilde\xi_t^R\bigr)=0,
\end{equation}
其中 $E_i^R:=E_i^R(u_t^R,\widetilde\xi_t^R)$ 如
注~\ref{rem:10-5-5}。再简记
\begin{equation}\label{eq:10-6-19}
  \widehat\xi_t^R
  :=\nabla_t\widetilde\xi_t^R
    -T_t^R\widetilde\nabla_t\eta_t^R.
\end{equation}
于是，由式~\eqref{eq:10-6-14}、\eqref{eq:10-6-15} 与
\eqref{eq:10-6-17}，有
\begin{equation}\label{eq:10-6-20}
  \|\widehat\xi_t^R\|_{1,p,R}
  \le c_2c_3c_5(\delta R)^{-2/p}
  \left(
    \|\partial_tu_t^0\|_{L^2}
    +\|\partial_tu_t^\infty\|_{L^2}
  \right).
\end{equation}
此外，把式~\eqref{eq:10-6-18} 与~\eqref{eq:10-6-19} 合并，得到
\begin{equation}\label{eq:10-6-21}
  \widetilde D_t^RE_2^RT_t^R\widetilde\nabla_t\eta_t^R
  =-\widetilde D_t^R
    \bigl(E_1^R\partial_tu_t^R+E_2^R\widehat\xi_t^R\bigr).
\end{equation}

下面断言，上式中出现的乘积 $\widetilde D E T$ 是 $DT$ 的小扰动，因而
可逆。为证明这一点，记
\[
  \Phi^R:=\Phi^R(u_t^R,\widetilde\xi_t^R):
  T_{u_t^R}M\longrightarrow T_{\widetilde u_t^R}M
\]
为关于\BookRecordIndex{T0052}埃尔米特联络 $\widetilde\nabla$ 沿测地线\BookRecordIndex{T0263}平行移动所确定的复向量丛
同构。于是算子
\[
  (\Phi^R)^{-1}\widetilde D_t^RE_i^R:
  W_{u_t^R}^{1,p}\longrightarrow L_{u_t^R,J^R}^p
\]
在算子范数意义下是 $D_t^R$ 的小扰动。更准确地说，由
引理~\ref{lem:10-3-1}，对 $i=1,2$ 有估计
\begin{equation}\label{eq:10-6-22}
  \bigl\|
    \widetilde D_t^RE_i^R(u_t^R,\widetilde\xi_t^R)\zeta
    -\Phi^R(u_t^R,\widetilde\xi_t^R)D_t^R\zeta
  \bigr\|_{0,p,R}
  \le c_6\|\widetilde\xi_t^R\|_{1,p,R}\|\zeta\|_{1,p,R}.
\end{equation}
将此式与命题~\ref{prop:10-5-1} 及式~\eqref{eq:10-6-14} 合并，得到
\[
\begin{aligned}
 &\|\widetilde D_t^RE_2^RT_t^R\eta-\Phi^R\eta\|_{0,p,R}\\
 &\quad\le
   \|\Phi^R(D_t^RT_t^R\eta-\eta)\|_{0,p,R}
   +\|\widetilde D_t^RE_2^RT_t^R\eta
      -\Phi^RD_t^RT_t^R\eta\|_{0,p,R}\\
 &\quad\le\frac12\|\eta\|_{0,p,R}
   +c_6\|\widetilde\xi_t^R\|_{1,p,R}
          \|T_t^R\eta\|_{1,p,R}\\
 &\quad\le\frac12
   \bigl(1+c_0c_2c_6(\delta R)^{-2/p}\bigr)
   \|\eta\|_{0,p,R}.
\end{aligned}
\]
若 $c_0c_2c_6(\delta R)^{-2/p}\le1/2$，则算子
\[
  \widetilde D_t^RE_2^RT_t^R:
  L_{u_t^R,J^R}^p\longrightarrow L_{\widetilde u_t^R,J^R}^p
\]
可逆，且其逆的范数不超过 $4$。因此，由式~\eqref{eq:10-6-21}，
\[
\begin{aligned}
  \|\widetilde\nabla_t\eta_t^R\|_{0,p,R}
  \le{}&4\left(
    \|\widetilde D_t^R(E_1^R\partial_tu_t^R)\|_{0,p,R}
    +\|\widetilde D_t^R(E_2^R\widehat\xi_t^R)\|_{0,p,R}
  \right)\\
  \le{}&4\|\Phi^RD_t^R\partial_tu_t^R\|_{0,p,R}
    +4c_6\|\widetilde\xi_t^R\|_{1,p,R}
            \|\partial_tu_t^R\|_{1,p,R}
    +c_7\|\widehat\xi_t^R\|_{1,p,R}\\
  \le{}&4c_4(\delta R)^{-2/p}\|\partial_tu_t^R\|_{1,\infty,R}
    +4c_2c_6(\delta R)^{-2/p}\|\partial_tu_t^R\|_{1,p,R}\\
  &\quad+c_2c_3c_5c_7(\delta R)^{-2/p}
    \left(
      \|\partial_tu_t^0\|_{L^2}
      +\|\partial_tu_t^\infty\|_{L^2}
    \right)\\
  \le{}&c_8(\delta R)^{-2/p}
    \left(
      \|\partial_tu_t^0\|_{L^2}
      +\|\partial_tu_t^\infty\|_{L^2}
    \right).
\end{aligned}
\]
这里第二个不等式用到了式~\eqref{eq:10-6-22} 以及
习题~\ref{ex:10-3-3} 所给 $\widetilde D_t^R$ 的一致有界性；第三个不等式
由式~\eqref{eq:10-6-14}、\eqref{eq:10-6-16}、\eqref{eq:10-6-20} \BookRecordIndex{T0326}推出；
最后一个不等式则由式~\eqref{eq:10-6-15} \BookRecordIndex{T0326}推出。将它与
式~\eqref{eq:10-6-19} 合并，得到
\[
\begin{aligned}
  \|\nabla_t\widetilde\xi_t^R\|_{1,p,R}
  &\le\|\widehat\xi_t^R\|_{1,p,R}
     +\frac{c_0}{2}\|\widetilde\nabla_t\eta_t^R\|_{0,p,R}\\
  &\le c_9(\delta R)^{-2/p}
    \left(
      \|\partial_tu_t^0\|_{L^2}
      +\|\partial_tu_t^\infty\|_{L^2}
    \right).
\end{aligned}
\]
这里还用到了式~\eqref{eq:10-6-20}。这就证明了
命题~\ref{prop:10-5-4}。
\end{proof}

\begin{proof}[定理~\ref{thm:10-1-2} (iv) 的证明]
最后简述\BookRecordIndex{T0254}粘合嵌入 $\iota_c^R$ 保持定向的证明。这里再次把 $\iota_c^R$
分解为\BookRecordIndex{T0394}预粘合映射 $f^R$ 与近似映射；在
$u=(u^0,u^\infty)\in\mathcal M(c)$ 处，其导数
\[
  d\iota_c^R(u):\ker D_u\longrightarrow\ker D_{\widetilde u^R}
\]
相应地分解为引理~\ref{lem:10-6-3} 中的同构
\begin{equation}\label{eq:10-6-23}
  (\operatorname{id}-Q_{u^R}D_{u^R})df^R(u):
  \ker D_u\longrightarrow\ker D_{u^R},
\end{equation}
其中 $u^R:=f^R(u)$；以及一个同构
\begin{equation}\label{eq:10-6-24}
  \ker D_{u^R}\longrightarrow\ker D_{\widetilde u^R}.
\end{equation}
命题~\ref{prop:10-5-4} 证明中的估计表明，在核丛的一个适当平凡化中，
同构~\eqref{eq:10-6-24} 接近恒等映射，因而保持定向。另一方面，由
习题~\ref{ex:10-6-4}，\BookRecordIndex{T0394}预粘合同构~\eqref{eq:10-6-23} 也保持定向。
所以，只要 $\delta R$ 充分大，$\iota_c^R$ 就是保持定向的嵌入。这证明了
定理~\ref{thm:10-1-2} (iv)。
\end{proof}
